\section{Narrativa e \textit{storytelling}}
\label{sec:narrativa}

A orientação acolhe o \textit{storytelling}. O grupo adota a estrutura clássica de situação, conflito e resolução, defendida por \citeonline{duarte2010} como forma de conduzir o público de um estado atual a um estado desejado. A história é protagonizada por uma tutora fictícia, \textbf{Maria}, e seu cão, \textbf{Thor}. Trata-se de uma persona composta, construída a partir dos relatos das entrevistas e dos testes da Atividade VI, e não de uma pessoa real; a lâmina 2 informa esse caráter, para evitar que a narrativa seja confundida com um caso documentado.

\begin{description}
  \item[Situação.] Maria mora em bairro de baixa renda da Região Metropolitana de Belo Horizonte e ama o Thor. Como a consulta particular custa de R\$~150 a R\$~350 na capital, as vacinas e as consultas de rotina vão sendo adiadas.
  \item[Conflito.] Thor adoece. Sem orientação, Maria não sabe se é grave; sem dinheiro, não sabe aonde ir. O que poderia ter sido prevenido vira emergência.
  \item[Resolução.] Com o PetCuida, Maria recebe lembretes de vacina, descreve o sintoma na triagem orientativa, é encaminhada a uma clínica parceira, paga o preço social e usa créditos solidários. A história termina antes da emergência, e não depois dela.
\end{description}

O arco retorna ao final: a lâmina de equipe e pedido retoma a imagem de Maria e do Thor, o que dá fechamento à apresentação. A narrativa mantém a honestidade do produto: o aplicativo orienta e encaminha, mas não diagnostica nem substitui o médico-veterinário, aviso que aparece na lâmina da solução.
